%!TeX root=../wowtop.tex

\ArtChapter[Under Foot]{18head}

\lettrine[lines=4,findent=2pt]{I}{n} the first book I have wandered so much from my own adventures to tell of the experiences of my brother that all through the last two chapters I and the curate have been lurking in the empty house at Halliford whither we fled to escape the Black Smoke. There I will resume. We stopped there all Sunday night and all the next day—the day of the panic—in a little island of daylight, cut off by the Black Smoke from the rest of the world. We could do nothing but wait in aching inactivity during those two weary days.

My mind was occupied by anxiety for my wife. I figured her at Leatherhead, terrified, in danger, mourning me already as a dead man. I paced the rooms and cried aloud when I thought of how I was cut off from her, of all that might happen to her in my absence. My cousin I knew was brave enough for any emergency, but he was not the sort of man to realise danger quickly, to rise promptly. What was needed now was not bravery, but circumspection. My only consolation was to believe that the Martians were moving Londonward and away from her. Such vague anxieties keep the mind sensitive and painful. I grew very weary and irritable with the curate's perpetual ejaculations; I tired of the sight of his selfish despair. After some ineffectual remonstrance I kept away from him, staying in a room—evidently a children's schoolroom—containing globes, forms, and copybooks. When he followed me thither, I went to a box room at the top of the house and, in order to be alone with my aching miseries, locked myself in.

We were hopelessly hemmed in by the Black Smoke all that day and the morning of the next. There were signs of people in the next house on Sunday evening—a face at a window and moving lights, and later the slamming of a door. But I do not know who these people were, nor what became of them. We saw nothing of them next day. The Black Smoke drifted slowly riverward all through Monday morning, creeping nearer and nearer to us, driving at last along the roadway outside the house that hid us.

A Martian came across the fields about midday, laying the stuff with a jet of superheated steam that hissed against the walls, smashed all the windows it touched, and scalded the curate's hand as he fled out of the front room. When at last we crept across the sodden rooms and looked out again, the country northward was as though a black snowstorm had passed over it. Looking towards the river, we were astonished to see an unaccountable redness mingling with the black of the scorched meadows.

For a time we did not see how this change affected our position, save that we were relieved of our fear of the Black Smoke. But later I perceived that we were no longer hemmed in, that now we might get away. So soon as I realised that the way of escape was open, my dream of action returned. But the curate was lethargic, unreasonable.

»We are safe here,« he repeated; »safe here.«

I resolved to leave him—would that I had! Wiser now for the artilleryman's teaching, I sought out food and drink. I had found oil and rags for my burns, and I also took a hat and a flannel shirt that I found in one of the bedrooms. When it was clear to him that I meant to go alone—had reconciled myself to going alone—he suddenly roused himself to come. And all being quiet throughout the afternoon, we started about five o'clock, as I should judge, along the blackened road to Sunbury.

In Sunbury, and at intervals along the road, were dead bodies lying in contorted attitudes, horses as well as men, overturned carts and luggage, all covered thickly with black dust. That pall of cindery powder made me think of what I had read of the destruction of Pompeii. We got to Hampton Court without misadventure, our minds full of strange and unfamiliar appearances, and at Hampton Court our eyes were relieved to find a patch of green that had escaped the suffocating drift. We went through Bushey Park, with its deer going to and fro under the chestnuts, and some men and women hurrying in the distance towards Hampton, and so we came to Twickenham. These were the first people we saw.

Away across the road the woods beyond Ham and Petersham were still afire. Twickenham was uninjured by either Heat-Ray or Black Smoke, and there were more people about here, though none could give us news. For the most part they were like ourselves, taking advantage of a lull to shift their quarters. I have an impression that many of the houses here were still occupied by scared inhabitants, too frightened even for flight. Here too the evidence of a hasty rout was abundant along the road. I remember most vividly three smashed bicycles in a heap, pounded into the road by the wheels of subsequent carts. We crossed Richmond Bridge about half past eight. We hurried across the exposed bridge, of course, but I noticed floating down the stream a number of red masses, some many feet across. I did not know what these were—there was no time for scrutiny—and I put a more horrible interpretation on them than they deserved. Here again on the Surrey side were black dust that had once been smoke, and dead bodies—a heap near the approach to the station; but we had no glimpse of the Martians until we were some way towards Barnes.

We saw in the blackened distance a group of three people running down a side street towards the river, but otherwise it seemed deserted. Up the hill Richmond town was burning briskly; outside the town of Richmond there was no trace of the Black Smoke.

Then suddenly, as we approached Kew, came a number of people running, and the upperworks of a Martian fighting-machine loomed in sight over the housetops, not a hundred yards away from us. We stood aghast at our danger, and had the Martian looked down we must immediately have perished. We were so terrified that we dared not go on, but turned aside and hid in a shed in a garden. There the curate crouched, weeping silently, and refusing to stir again.

But my fixed idea of reaching Leatherhead would not let me rest, and in the twilight I ventured out again. I went through a shrubbery, and along a passage beside a big house standing in its own grounds, and so emerged upon the road towards Kew. The curate I left in the shed, but he came hurrying after me.

\begin{wrapfigure}{O}{0.5\textwidth}
\centering
\includegraphics[width=0.5\textwidth]{18tripod}
\end{wrapfigure}

That second start was the most foolhardy thing I ever did. For it was manifest the Martians were about us. No sooner had the curate overtaken me than we saw either the fighting-machine we had seen before or another, far away across the meadows in the direction of Kew Lodge. Four or five little black figures hurried before it across the green-grey of the field, and in a moment it was evident this Martian pursued them. In three strides he was among them, and they ran radiating from his feet in all directions. He used no Heat-Ray to destroy them, but picked them up one by one. Apparently he tossed them into the great metallic carrier which projected behind him, much as a workman's basket hangs over his shoulder.

It was the first time I realised that the Martians might have any other purpose than destruction with defeated humanity. We stood for a moment petrified, then turned and fled through a gate behind us into a walled garden, fell into, rather than found, a fortunate ditch, and lay there, scarce daring to whisper to each other until the stars were out.

I suppose it was nearly eleven o'clock before we gathered courage to start again, no longer venturing into the road, but sneaking along hedgerows and through plantations, and watching keenly through the darkness, he on the right and I on the left, for the Martians, who seemed to be all about us. In one place we blundered upon a scorched and blackened area, now cooling and ashen, and a number of scattered dead bodies of men, burned horribly about the heads and trunks but with their legs and boots mostly intact; and of dead horses, fifty feet, perhaps, behind a line of four ripped guns and smashed gun carriages.

Sheen, it seemed, had escaped destruction, but the place was silent and deserted. Here we happened on no dead, though the night was too dark for us to see into the side roads of the place. In Sheen my companion suddenly complained of faintness and thirst, and we decided to try one of the houses.

The first house we entered, after a little difficulty with the window, was a small semi-detached villa, and I found nothing eatable left in the place but some mouldy cheese. There was, however, water to drink; and I took a hatchet, which promised to be useful in our next house-breaking.

We then crossed to a place where the road turns towards Mortlake. Here there stood a white house within a walled garden, and in the pantry of this domicile we found a store of food—two loaves of bread in a pan, an uncooked steak, and the half of a ham. I give this catalogue so precisely because, as it happened, we were destined to subsist upon this store for the next fortnight. Bottled beer stood under a shelf, and there were two bags of haricot beans and some limp lettuces. This pantry opened into a kind of wash-up kitchen, and in this was firewood; there was also a cupboard, in which we found nearly a dozen of burgundy, tinned soups and salmon, and two tins of biscuits.

We sat in the adjacent kitchen in the dark—for we dared not strike a light—and ate bread and ham, and drank beer out of the same bottle. The curate, who was still timorous and restless, was now, oddly enough, for pushing on, and I was urging him to keep up his strength by eating when the thing happened that was to imprison us.

»It can't be midnight yet,« I said, and then came a blinding glare of vivid green light. Everything in the kitchen leaped out, clearly visible in green and black, and vanished again. And then followed such a concussion as I have never heard before or since. So close on the heels of this as to seem instantaneous came a thud behind me, a clash of glass, a crash and rattle of falling masonry all about us, and the plaster of the ceiling came down upon us, smashing into a multitude of fragments upon our heads. I was knocked headlong across the floor against the oven handle and stunned. I was insensible for a long time, the curate told me, and when I came to we were in darkness again, and he, with a face wet, as I found afterwards, with blood from a cut forehead, was dabbing water over me.

\begin{figure}[tb!]
\centering
\includegraphics[width=\textwidth]{18kaboom}
\end{figure}

For some time I could not recollect what had happened. Then things came to me slowly. A bruise on my temple asserted itself.

»Are you better?« asked the curate in a whisper.

At last I answered him. I sat up.

»Don't move,« he said. »The floor is covered with smashed crockery from the dresser. You can't possibly move without making a noise, and I fancy \textit{they} are outside.«

We both sat quite silent, so that we could scarcely hear each other breathing. Everything seemed deadly still, but once something near us, some plaster or broken brickwork, slid down with a rumbling sound. Outside and very near was an intermittent, metallic rattle.

»That!« said the curate, when presently it happened again.

»Yes,« I said. »But what is it?«

»A Martian!« said the curate.

I listened again.

»It was not like the Heat-Ray,« I said, and for a time I was inclined to think one of the great fighting-machines had stumbled against the house, as I had seen one stumble against the tower of Shepperton Church.

Our situation was so strange and incomprehensible that for three or four hours, until the dawn came, we scarcely moved. And then the light filtered in, not through the window, which remained black, but through a triangular aperture between a beam and a heap of broken bricks in the wall behind us. The interior of the kitchen we now saw greyly for the first time.

The window had been burst in by a mass of garden mould, which flowed over the table upon which we had been sitting and lay about our feet. Outside, the soil was banked high against the house. At the top of the window frame we could see an uprooted drainpipe. The floor was littered with smashed hardware; the end of the kitchen towards the house was broken into, and since the daylight shone in there, it was evident the greater part of the house had collapsed. Contrasting vividly with this ruin was the neat dresser, stained in the fashion, pale green, and with a number of copper and tin vessels below it, the wallpaper imitating blue and white tiles, and a couple of coloured supplements fluttering from the walls above the kitchen range.

As the dawn grew clearer, we saw through the gap in the wall the body of a Martian, standing sentinel, I suppose, over the still glowing cylinder. At the sight of that we crawled as circumspectly as possible out of the twilight of the kitchen into the darkness of the scullery.

Abruptly the right interpretation dawned upon my mind.

»The fifth cylinder,« I whispered, »the fifth shot from Mars, has struck this house and buried us under the ruins!«

For a time the curate was silent, and then he whispered:

»God have mercy upon us!«

I heard him presently whimpering to himself.

Save for that sound we lay quite still in the scullery; I for my part scarce dared breathe, and sat with my eyes fixed on the faint light of the kitchen door. I could just see the curate's face, a dim, oval shape, and his collar and cuffs. Outside there began a metallic hammering, then a violent hooting, and then again, after a quiet interval, a hissing like the hissing of an engine. These noises, for the most part problematical, continued intermittently, and seemed if anything to increase in number as time wore on. Presently a measured thudding and a vibration that made everything about us quiver and the vessels in the pantry ring and shift, began and continued. Once the light was eclipsed, and the ghostly kitchen doorway became absolutely dark. For many hours we must have crouched there, silent and shivering, until our tired attention failed\textellipsis

At last I found myself awake and very hungry. I am inclined to believe we must have spent the greater portion of a day before that awakening. My hunger was at a stride so insistent that it moved me to action. I told the curate I was going to seek food, and felt my way towards the pantry. He made me no answer, but so soon as I began eating the faint noise I made stirred him up and I heard him crawling after me.

\begin{figure}[b!]
\centering
\includegraphics[width=.8\linewidth]{18tailpiece}
%\captionlistentry{Tailpiece to Chapter \thechapter}
\end{figure}